% paper_mu_relation.tex
% The Khronon Mass Scale from Recombination
% Author: Sheng-Kai Huang
% Compile: pdflatex paper_mu_relation.tex

\documentclass[prl,twocolumn,superscriptaddress,longbibliography]{revtex4-2}

\usepackage{amsmath}
\usepackage{amssymb}
\usepackage{amsfonts}
\usepackage{bm}
\usepackage{hyperref}
\usepackage{booktabs}

% Macros
\newcommand{\kB}{k_{\mathrm{B}}}
\newcommand{\Tr}{\mathrm{Tr}}
\newcommand{\Mpc}{\,\mathrm{Mpc}}
\newcommand{\zdec}{z_{\mathrm{dec}}}
\newcommand{\rd}{r_{\mathrm{d}}}
\newcommand{\Ob}{\Omega_{\mathrm{b}}}
\newcommand{\Odm}{\Omega_{\mathrm{DM}}}

\begin{document}

\title{The Khronon Mass Scale from Recombination:\\
Two Numerical Relations Connecting MOND, CMB,
and Dark Matter}

\author{Sheng-Kai Huang}
\email{akai@fawstudio.com}
\affiliation{Independent Researcher}

\date{\today}

\begin{abstract}
In the Blanchet--Skordis Khronon dark matter theory,
the inverse mass parameter $\mu^{-1} = 22.3\Mpc$ is
fitted to galaxy rotation curves and CMB data, but
its theoretical origin is unknown.
We report two numerical relations, apparently absent
from the literature, that express $\mu$ in terms of
fundamental cosmological quantities:
(i)~$\mu c^2 = a_0\,(1 + \zdec)$,
relating the Khronon mass to the MOND acceleration
blueshifted to photon decoupling ($\mu^{-1} = 22.25\Mpc$,
$0.2\%$ agreement); and
(ii)~$\mu = 2\pi(1 + \Ob)/\rd$,
relating the Khronon Compton wavelength to the
baryon-loaded sound horizon
($\mu^{-1} = 22.31\Mpc$, $0.04\%$ agreement).
The two relations are structurally independent yet
mutually consistent to $0.27\%$.
Both connect a galactic-scale parameter ($\mu$ or $a_0$)
to recombination-epoch quantities ($\zdec$, $\rd$, $\Ob$),
suggesting that the Khronon mass is not a free parameter
but is fixed by the physics of baryon--photon decoupling.
\end{abstract}

\maketitle


%=======================================================================
\section{Introduction}
\label{sec:intro}
%=======================================================================

The dark matter problem admits two complementary
phenomenological descriptions.
On galactic scales, the MOND acceleration
$a_0 \approx 1.2 \times 10^{-10}\;\mathrm{m\,s}^{-2}$
governs the baryonic Tully--Fisher relation and the
radial acceleration relation~\cite{Milgrom1983,McGaugh2016}.
On cosmological scales, the $\Lambda$CDM paradigm
reproduces the CMB power spectrum and large-scale structure
with a pressureless cold dark matter component,
$\Odm h^2 = 0.120$~\cite{Planck2018}.
Bridging these regimes has proven difficult:
theories that reproduce MOND phenomenology generically
fail CMB constraints, and vice versa.

A notable exception is the Blanchet--Skordis (BS) Khronon
theory~\cite{BS2024,BS2025}, which extends the
Skordis--Z\l o\'snik action~\cite{SZ2021} with a
Dirac--Born--Infeld (DBI) kinetic structure.
The theory contains two sectors: a galactic sector
$J(Y)$ that recovers MOND, and a cosmological sector
$K(Q) = \mu^2(Q - 1)^2$ that provides CDM-like
perturbations through ghost condensation.
The single parameter $\mu$ controls the transition
from the modified-gravity regime to the Newtonian regime;
its inverse $\mu^{-1} = 22.3\Mpc$ is determined by
fitting to CMB and rotation curve data~\cite{BS2025}.

Despite the theory's empirical success, the origin of $\mu$
remains unexplained.
Dimensional analysis suggests $\mu \sim H_0/c$, giving
$\mu^{-1} \sim c/H_0 \approx 4450\Mpc$---a factor of
${\sim}\,200$ too large, and incompatible with CMB
data~\cite{BS2025}.
In this Letter, we report two numerical relations
that pin down $\mu$ in terms of well-measured
cosmological quantities with sub-percent accuracy.


%=======================================================================
\section{Relation I: MOND acceleration at decoupling}
\label{sec:rel1}
%=======================================================================

We propose
\begin{equation}
  \mu\, c^2 \;=\; a_0\,(1 + \zdec) \,,
  \label{eq:rel1}
\end{equation}
where $\zdec = 1089.92$ is the photon decoupling
redshift~\cite{Planck2018}.
This can be read as: \emph{the Khronon mass scale equals
the MOND acceleration blueshifted to the epoch of
last scattering, expressed in natural units.}

Inverting Eq.~\eqref{eq:rel1} gives the Khronon
Compton wavelength,
\begin{equation}
  \mu^{-1}
  = \frac{c^2}{a_0\,(1 + \zdec)}
  = \frac{(2.998 \times 10^8)^2}
         {1.2 \times 10^{-10} \times 1090.92}
  = 22.25\Mpc .
  \label{eq:rel1_value}
\end{equation}
Compared to the BS2025 fitted value
$\mu^{-1}_{\mathrm{fit}} = 22.3\Mpc$~\cite{BS2025},
the agreement is
\begin{equation}
  \frac{|\mu^{-1} - \mu^{-1}_{\mathrm{fit}}|}
       {\mu^{-1}_{\mathrm{fit}}}
  = 0.2\% \,.
\end{equation}

\emph{Physical interpretation.}---The
quantity $a_0(1 + \zdec)$ is the
comoving MOND acceleration at decoupling.
In MOND phenomenology, $a_0$ marks
the transition between Newtonian and deep-MOND
dynamics at $z = 0$; at redshift $z$, this transition
scale blueshifts to $a_0(1 + z)$~\cite{Milgrom2020}.
Relation~\eqref{eq:rel1} states that the Khronon
mass $\mu$ is set by this transition scale evaluated
precisely at the epoch when baryons decouple from
photons.
This is physically natural: $\mu$ controls the
MOND--Newton transition in the Khronon theory,
and its value should be determined by the
cosmological epoch at which the relevant degrees
of freedom (baryonic perturbations) become
dynamically independent.


%=======================================================================
\section{Relation II: Baryon-loaded sound horizon}
\label{sec:rel2}
%=======================================================================

Independently, we find
\begin{equation}
  \mu \;=\; \frac{2\pi\,(1 + \Ob)}{\rd} \,,
  \label{eq:rel2}
\end{equation}
where $\rd = 147.09\Mpc$ is the comoving sound
horizon at the baryon drag epoch~\cite{Planck2018}
and $\Ob = 0.0493$ is the baryon density
parameter.  Inverting,
\begin{equation}
  \mu^{-1}
  = \frac{\rd}{2\pi\,(1 + \Ob)}
  = \frac{147.09}{2\pi \times 1.0493}
  = 22.31\Mpc .
  \label{eq:rel2_value}
\end{equation}
The agreement with the fitted value is
\begin{equation}
  \frac{|\mu^{-1} - \mu^{-1}_{\mathrm{fit}}|}
       {\mu^{-1}_{\mathrm{fit}}}
  = 0.04\% \,.
\end{equation}

\emph{Physical interpretation.}---The
sound horizon $\rd$ is the comoving distance
traveled by acoustic waves in the baryon--photon
plasma from the Big Bang to recombination.
Relation~\eqref{eq:rel2} states that the Khronon
Compton wavelength is $\rd / 2\pi$ corrected by
the baryon loading factor $(1 + \Ob)$.
In other words, $\mu^{-1}$ is the
\emph{reduced acoustic wavelength} of the sound
horizon, adjusted for the gravitational drag of
baryons on the plasma oscillations.
This connects the Khronon mass---which governs
galactic-scale MOND phenomenology---directly to the
BAO scale imprinted at recombination.


%=======================================================================
\section{Mutual consistency}
\label{sec:consistency}
%=======================================================================

Relations~\eqref{eq:rel1} and \eqref{eq:rel2}
are structurally independent:
the first involves $a_0$ and $\zdec$; the second
involves $\rd$ and $\Ob$.
Neither contains the other's characteristic quantities.
Yet the two predictions agree to
\begin{equation}
  \frac{|22.25 - 22.31|}{22.31} = 0.27\% \,,
\end{equation}
well within the observational uncertainty on $a_0$
(typically $5$--$10\%$~\cite{McGaugh2016,Lelli2017})
and on $\rd$ ($0.2\%$~\cite{Planck2018}).

Equating the two relations yields a
\emph{consistency condition}
linking MOND, CMB, and baryon physics:
\begin{equation}
  a_0
  = \frac{2\pi\,(1 + \Ob)\,c^2}
         {\rd\,(1 + \zdec)} \,.
  \label{eq:consistency}
\end{equation}
Evaluating with Planck~2018 parameters gives
\begin{equation}
  a_0^{\mathrm{pred}}
  = \frac{2\pi \times 1.0493 \times
          (2.998 \times 10^8)^2}
         {4.536 \times 10^{24} \times 1090.92}
  = 1.197 \times 10^{-10}\;\mathrm{m\,s}^{-2} ,
  \label{eq:a0_pred}
\end{equation}
where $\rd = 147.09\Mpc = 4.539 \times 10^{24}\;\mathrm{m}$.
This agrees with the empirical MOND value
$a_0 = (1.20 \pm 0.02_{\mathrm{stat}}
\pm 0.24_{\mathrm{sys}}) \times 10^{-10}\;\mathrm{m\,s}^{-2}$
\cite{McGaugh2016,Lelli2017} to $0.3\%$.

Equation~\eqref{eq:consistency} is a parameter-free
prediction: \emph{if} the Khronon mass is set by
recombination physics, \emph{then} the MOND
acceleration is calculable from $\rd$, $\zdec$,
$\Ob$, and $c$ alone.
This provides a sharp test---any future refinement
of $a_0$, $\rd$, or $\Ob$ must preserve
Eq.~\eqref{eq:consistency}.


%=======================================================================
\section{The factor-of-200 hierarchy}
\label{sec:hierarchy}
%=======================================================================

As noted in Sec.~\ref{sec:intro}, the naive dimensional
estimate $\mu \sim H_0/c$ overshoots by a factor
of ${\sim}\,200$.
Relation~\eqref{eq:rel1} makes this hierarchy transparent:
\begin{equation}
  \frac{\mu}{H_0/c}
  = \frac{a_0\,(1 + \zdec)}{c\, H_0}
  = \frac{(1 + \zdec)}{2\pi}
    \times \frac{a_0}{a_0^{\mathrm{M}}} \,,
  \label{eq:hierarchy}
\end{equation}
where $a_0^{\mathrm{M}} \equiv c\, H_0 / (2\pi)
\approx 1.04 \times 10^{-10}\;\mathrm{m\,s}^{-2}$ is
Milgrom's cosmological coincidence
value~\cite{Milgrom1983,Milgrom2020}.
Since $a_0 / a_0^{\mathrm{M}} \approx 1.15$ empirically,
\begin{equation}
  \frac{\mu}{H_0/c}
  \approx \frac{1091}{2\pi} \times 1.15
  \approx 200 \,.
  \label{eq:200}
\end{equation}
The factor of $200$ is thus not a mysterious large
number but the product of two well-understood
ingredients: $(1 + \zdec)/(2\pi)$, a purely
geometrical--cosmological ratio, and the
$O(1)$ correction $a_0/a_0^{\mathrm{M}}$.

This also clarifies why $\mu_0 = H_0/c$---the
value motivated by dimensional analysis and used
in Refs.~\cite{Huang2026_P3,Huang2026_P4}---is
excluded by CMB data: the physics of recombination
inserts a factor of $(1 + \zdec)/(2\pi)$ that
cannot be neglected.


%=======================================================================
\section{Discussion}
\label{sec:discussion}
%=======================================================================

\emph{Two routes to one number.}---The
two relations reported here arrive at
$\mu^{-1} \approx 22.3\Mpc$ via completely
different physical ingredients.
Relation~I uses the MOND acceleration $a_0$---a
galactic observable---and the decoupling redshift.
Relation~II uses the sound horizon $\rd$---a
cosmological observable---and the baryon fraction.
Their mutual consistency (Eq.~\eqref{eq:consistency})
is a nontrivial constraint that any microscopic
derivation of $\mu$ must satisfy.

\emph{Implications for the BS theory.}---If
either relation holds as an identity rather than a
coincidence, it reduces the free parameter count
of the BS Khronon theory from two ($a_0$, $\mu$) to
one ($a_0$ alone, with $\mu$ derived).
Furthermore, if the consistency condition
\eqref{eq:consistency} is exact,
$a_0$ itself is determined by cosmological
parameters, leaving the theory with
\emph{zero} free parameters in the dark sector
beyond $\Lambda$CDM inputs.

\emph{Connection to Milgrom's FUNDAMOND
program.}---Milgrom has long
advocated that $a_0 \sim c\, H_0$ reflects a deep
connection between local and cosmological
dynamics~\cite{Milgrom1999,Milgrom2020}.
Equation~\eqref{eq:consistency} refines
this conjecture: the connection runs not through
$H_0$ directly but through the sound horizon $\rd$,
the baryon fraction $\Ob$, and the decoupling
redshift $\zdec$.
All three are determined by the baryon--photon
plasma physics at recombination, suggesting that the
MOND scale $a_0$ is an \emph{imprint} of
recombination rather than a reflection of the
present-day Hubble rate.

\emph{What remains open.}---We have not derived
either relation from first principles.
A microscopic derivation would need to show that
ghost condensation in the Khronon sector locks onto
the acoustic scale at recombination.
One plausible mechanism is that the Khronon field,
being a timelike scalar coupled to the metric,
develops its condensate precisely when the
baryon--photon plasma decouples and the sound
speed drops sharply.
The coincidence $\mu^{-1} \sim \rd / (2\pi)$
suggests that the Khronon ``freezes in'' at a
wavelength corresponding to the last acoustic
oscillation, analogous to how axion dark matter
freezes at the QCD scale.
We leave a detailed derivation to future work.

\emph{Falsifiability.}---Both relations make
precise, testable predictions.
Relation~I predicts $\mu^{-1} = c^2 / [a_0(1 + \zdec)]$;
any independent measurement of $\mu$ from CMB or
large-scale structure data must match.
The consistency condition~\eqref{eq:consistency}
predicts $a_0 = 1.197 \times 10^{-10}\;\mathrm{m\,s}^{-2}$;
this is falsifiable by future high-precision measurements
of the radial acceleration relation (e.g., from DESI
or Euclid weak-lensing samples).


%=======================================================================
\section{Summary}
\label{sec:summary}
%=======================================================================

We have reported two numerical relations:
\begin{align}
  \text{(I)}\quad \mu\,c^2 &= a_0\,(1+\zdec) \,,
  \label{eq:summary1} \\
  \text{(II)}\quad \mu &= 2\pi\,(1+\Ob)/\rd \,,
  \label{eq:summary2}
\end{align}
each reproducing the BS-fitted Khronon mass
$\mu^{-1} = 22.3\Mpc$ to better than $0.3\%$.
Their mutual consistency yields a parameter-free
prediction for the MOND acceleration,
$a_0 = 2\pi(1+\Ob)\,c^2 / [\rd(1+\zdec)]
= 1.20 \times 10^{-10}\;\mathrm{m\,s}^{-2}$.
If confirmed as exact relations, they eliminate the last
free parameter of the Khronon dark matter theory and
anchor the MOND scale to the physics of
baryon--photon decoupling.


%=======================================================================
\begin{acknowledgments}
The author thanks Luc Blanchet and Constantinos Skordis
for making the BS Khronon theory publicly available
through detailed publications.
\end{acknowledgments}
%=======================================================================


\begin{thebibliography}{99}

\bibitem{Milgrom1983}
M.~Milgrom,
``A modification of the Newtonian dynamics as a possible
alternative to the hidden mass hypothesis,''
Astrophys.\ J.\ \textbf{270}, 365 (1983).

\bibitem{McGaugh2016}
S.~S.~McGaugh, F.~Lelli, and J.~M.~Schombert,
``Radial acceleration relation in rotationally supported
galaxies,''
Phys.\ Rev.\ Lett.\ \textbf{117}, 201101 (2016).

\bibitem{Lelli2017}
F.~Lelli, S.~S.~McGaugh, J.~M.~Schombert, and
M.~S.~Pawlowski,
``One law to rule them all: The radial acceleration
relation of galaxies,''
Astrophys.\ J.\ \textbf{836}, 152 (2017).

\bibitem{Planck2018}
N.~Aghanim \textit{et al.} (Planck Collaboration),
``Planck 2018 results. VI. Cosmological parameters,''
Astron.\ Astrophys.\ \textbf{641}, A6 (2020);
arXiv:1807.06209.

\bibitem{SZ2021}
C.~Skordis and T.~Z\l o\'snik,
``New relativistic theory for modified Newtonian
dynamics,''
Phys.\ Rev.\ Lett.\ \textbf{127}, 161302 (2021).

\bibitem{BS2024}
L.~Blanchet and C.~Skordis,
``Dark matter as a Khronon condensate,''
arXiv:2404.06584 (2024).

\bibitem{BS2025}
L.~Blanchet and C.~Skordis,
``Khronon dark matter: CMB and large-scale structure,''
arXiv:2507.00912 (2025).

\bibitem{Milgrom1999}
M.~Milgrom,
``The modified dynamics as a vacuum effect,''
Phys.\ Lett.\ A \textbf{253}, 273 (1999).

\bibitem{Milgrom2020}
M.~Milgrom,
``MOND vs.\ dark matter in light of historical parallels,''
Stud.\ Hist.\ Philos.\ Mod.\ Phys.\ \textbf{71}, 170 (2020);
arXiv:1910.04368.

\bibitem{Huang2026_P3}
S.-K.~Huang,
``Dark matter as Khronon condensate: Weak-field
phenomenology of the $\tau$ framework,''
in preparation (2026).

\bibitem{Huang2026_P4}
S.-K.~Huang,
``The cosmological sector of the Khronon $\tau$ framework:
$\mu_0 = H_0/c$ and dark matter density,''
in preparation (2026).

\end{thebibliography}

\end{document}
